\begin{question}

A spin $\frac 1 2$ system is known to be in an eigenstate of
$\vectorbold{S}\cdot\vectorbold{\hat{n}}$ with eigenvalue $\hbar/2$,
where $\vectorbold{\hat{n}}$ is a unit vector lying in the xz-plane
that makes an angle $\gamma$ with the positive z-axis.

\begin{questionpart}
  Suppose $S_x$ is measured. What is the probability of getting $+\hbar/2$?
\end{questionpart}

\begin{questionpart}
  Evaluate the dispersion in $S_x$, that is
  \begin{equation*}
    \expval{(S_x - \expval{S_x})^2}
  \end{equation*}
\end{questionpart}

(For your own peace of mind check your answers for the special cases
$\gamma = 0\text{, }\pi/2\text{, and }\pi.$)
\end{question}

\begin{purpose}
  Learn how to apply our result from problem \ref{1.11:theQuestion}.
\end{purpose}

\begin{answer}

  \begin{answerpart}{Probability}
    Let's start by putting down a few equations we get from the
    problem definition.  We know that $\vectorbold{\hat{n}}$ \emph{is}
    in the $xz$ plane, which means that it is \emph{not} at all on the
    $y$ axis.

    \begin{equation}
      \vectorbold{\hat{n}} \cdot \vectorbold{\hat{y}} = 0
    \end{equation}

    \begin{equation}
      || \vectorbold{\hat{n}} \cdot \vectorbold{\hat{x}}||^2
      + || \vectorbold{\hat{n}} \cdot \vectorbold{\hat{z}}||^2
      = 1
    \end{equation}

    Since we're using $S_x$ as our measurement basis, let's pull up
    those definitions also:

    \begin{equation*}
      \tag{1.110a}
      \begin{split}
        \ket{x\pm} &= \frac{1}{\sqrt{2}} [\ket + \pm \ket -]\\
        &=> \ket{x+} = \frac{1}{\sqrt{2}} [\ket + + \ket -]
      \end{split}
    \end{equation*}

    \begin{equation*}
      \tag{1.111a}
      S_x = \frac{\hbar}{2} [\ket - \bra + + \ket + \bra - ]
    \end{equation*}

    Referencing our axis diagram (\qref{spinAxes}) in the QRF, we see
    that the problem describes an azimuth angle ($\alpha$) of $0$ and
    an elevation-complement angle ($\beta$) of $\gamma$.  Applying our
    result from problem \ref{1.11:theQuestion}
    (\ref{1.11:spinAngle}/\qref{spinAngle}):

    \begin{equation}
      \begin{split}
        \ket{\alpha_+}
        &= \cos{\frac{\beta}{2}} \ket + + e^{-i\alpha}\sin{\frac{\beta}{2}} \ket -\\
        &= \cos{\frac{\gamma}{2}} \ket + + \sin{\frac{\gamma}{2}} \ket -\\
      \end{split}
    \end{equation}

    Evaluating our probability:
    \begin{equation}
      \begin{split}
        P(x+)
        &= \left|\braket{x+}{\alpha+}\right|^2\\
        &= \left|\frac{1}{\sqrt{2}}
        (\bra + + \bra -)
        (\cos{\frac{\gamma}{2}} \ket + + \sin{\frac{\gamma}{2}} \ket -)\right|^2\\
        &= \frac{1}{2}\left|
        \cos{\frac{\gamma}{2}}\braket +
        + \sin{\frac{\gamma}{2}}\braket -\right|^2\\
        &= \frac{1}{2}\left|
        \cos{\frac{\gamma}{2}} + \sin{\frac{\gamma}{2}}
        \right|^2\\
        &= \frac{1}{2}\left|
        (\cos^2{\frac{\gamma}{2}} + \sin^2{\frac{\gamma}{2}})
        + 2\cos{\frac{\gamma}{2}}\sin{\frac{\gamma}{2}}
        \right|\\
        &= \frac{1}{2}(1 + \sin{\gamma})\\
      \end{split}
    \end{equation}

    Obviously we can evaluate the general solution, but what fun is
    that?  

    \begin{equation}
      \begin{alignedat}{2}
      %
      &&\braket{x+}{\alpha_+} &=
      \bra{x+}\left(\cos{\frac{\gamma}{2}} \ket +
      + \sin{\frac{\gamma}{2}} \ket -\right)\\
      %
      \gamma &= 0 \quad &
      \braket{x+}{\alpha_+} &=
      \bra{x+}(\cos{0} \ket + + \sin{0} \ket -)\\
      & & &=
      \frac{1}{\sqrt{2}} (\bra + + \bra -)(\cos{0} \ket + + \sin{0} \ket -)\\
      %
      & & &=
      \frac{1}{\sqrt{2}}\\
      %
      \gamma &= \frac{\pi}{2} \quad &
      \braket{x+}{\alpha_+} &=
      \bra{x+}(\cos{\frac{\pi}{4}} \ket + + \sin{\frac{\pi}{4}} \ket -)\\
      & & &=
      \frac{1}{\sqrt{2}}(\bra + + \bra -)\frac{1}{\sqrt{2}}(\ket + + \ket -)\\
      & & &=
      \frac{1}{2}(\braket + + \braket -)\\
      & & &=
      1\\
      %
      \gamma &= \pi \quad &
      \braket{x+}{\alpha_+} &=
      \bra{x+}\left(\cos{\frac{\pi}{2}} \ket +
      + \sin{\frac{\pi}{2}} \ket -\right)\\
      & & &=
      \frac{1}{\sqrt{2}}(\bra + + \bra -)\left(\cos{\frac{\pi}{2}} \ket +
      + \sin{\frac{\pi}{2}} \ket -\right)\\
      & & &=
      \frac{1}{\sqrt{2}}\\
      %
    \end{alignedat}
  \end{equation}

  Translating those magnitudes into probabilities, we see that when we
  are aligned with $x+$, $P(\frac{\hbar}{2}) = 1$ whereas when we are
  perpendicular to it, our probability is $\frac 1 2$.  When we're
  perpendicular, then we have no information about which way it'll go
  (either $\frac \hbar 2$ or $- \frac \hbar 2$) so we expect a
  probability of $\frac 1 2$.

  \begin{equation}
    \begin{alignedat}{2}
      P(x+;\gamma&=0) & &= \frac 1 2\\
      P(x+;\gamma&=\frac{\pi}{2}) & &= 1\\
      P(x+;\gamma&=\pi) & &= \frac 1 2\\
    \end{alignedat}
  \end{equation}

  \end{answerpart}

  \begin{answerpart}{Dispersion}
  \end{answerpart}

  
\end{answer}
